\documentclass[11pt,a4paper]{article}
\usepackage[T1]{fontenc}
\usepackage[utf8]{inputenc}
\usepackage{lmodern}
\usepackage[margin=27mm,headheight=14pt]{geometry}
\usepackage{amsmath,amssymb,amsthm,mathtools}
\usepackage{microtype,booktabs,tabularx,longtable,array}
\usepackage{xcolor,enumitem,fancyhdr,hyperref}
\hypersetup{colorlinks=true,linkcolor=blue!45!black,citecolor=blue!45!black,
 urlcolor=blue!45!black,
 pdftitle={Auxiliary Problems and Arithmetic Barriers},
 pdfauthor={Research note prepared for Keio},
 pdfsubject={Sharp gap-mass threshold, sparse exception series, reciprocal tails, and a dated priority audit}}
\urlstyle{same}
\pagestyle{fancy}\fancyhf{}
\fancyhead[L]{\small Auxiliary problems and arithmetic barriers}
\fancyhead[R]{\small Research note}
\fancyfoot[C]{\thepage}
\renewcommand{\headrulewidth}{0.3pt}
\setlength{\emergencystretch}{3em}
\setlist{itemsep=3pt,topsep=5pt}
\newtheorem{theorem}{Theorem}[section]
\newtheorem{proposition}[theorem]{Proposition}
\newtheorem{lemma}[theorem]{Lemma}
\newtheorem{corollary}[theorem]{Corollary}
\theoremstyle{definition}
\newtheorem{definition}[theorem]{Definition}
\newtheorem{question}[theorem]{Auxiliary question}
\theoremstyle{remark}\newtheorem{remark}[theorem]{Remark}
\newcommand{\N}{\mathbb N}\newcommand{\Z}{\mathbb Z}
\newcommand{\Q}{\mathbb Q}\newcommand{\C}{\mathbb C}
\newcommand{\E}{\mathcal E}
\newcommand{\ord}{\operatorname{ord}}
\newcommand{\floor}[1]{\left\lfloor #1\right\rfloor}
\newcommand{\abs}[1]{\left|#1\right|}
\newcommand{\statusbox}[1]{\begin{center}\fcolorbox{blue!45!black}{blue!3}{\begin{minipage}{0.91\linewidth}\small #1\end{minipage}}\end{center}}

\title{\textbf{Auxiliary Problems and\\Arithmetic Barriers}\\[7pt]
\large A sharp gap-mass threshold, sparse exception series,\\
rational reciprocal tails, and a priority audit}
\author{Research note prepared for Keio}
\date{27 November 2026}

\begin{document}
\maketitle\thispagestyle{empty}
\begin{abstract}
We investigate whether auxiliary problems suggested by the Verifiable
Mathematical Reasoning Corpus can close an existing named conjecture.
For an infinite integer set $Q$, let $S_Q(x)$ be the total length of its
consecutive gaps of size at least $\sqrt x$ whose left endpoints lie in
$[x,2x]$. We prove that infinitely many empty consecutive-square
intervals force $\limsup S_Q(x)/\sqrt x\ge2\sqrt2$. This constant is
sharp even among sets with weighted counting function
$x+O(\sqrt x)$, counting function $\int_3^xdt/\log t+
O(\sqrt x/\log x)$, and any prescribed finite prime prefix.
For the actual Legendre exception series, we prove that finite support,
polynomiality, rationality, algebraicity, and differential finiteness
are equivalent, and give quantitative finite-certificate criteria.
Explicit thin bases and sparse words disprove two proposed routes from
sparsity alone. In a second arithmetic direction, we reconstruct the
bounded-negative-defect theorem for rational reciprocal tails, with an
elementary prime-support argument and full attribution to earlier work.
Exact graph calculations provide additional tests, but the corresponding
named conjectures have prior resolutions or public full-proof claims.
No new unconditional resolution satisfying all six requested conditions
G1--G6 is obtained. The paper distinguishes complete auxiliary proofs
from missing transfers to the original problems.
\end{abstract}

\statusbox{\textbf{Research outcome.}
The sharp constant $2\sqrt2$ is established for an auxiliary extremal
problem about general integer sets, not as a new theorem about the
distribution of the primes. Neither Legendre's conjecture nor the
unrestricted Erd\H{o}s--Graham reciprocal-sum problem is resolved here.
Results rediscovered during the investigation are attributed to their
identified predecessors. No claim of historical priority is made.}

\vfill
\noindent\small\textbf{Reproducibility.}
This is a self-contained English manuscript with an embedded bibliography.
The accompanying single \texttt{.tex} file requires no external figures,
data files, or bibliography database. All arguments supplied here are
mathematical proofs; no numerical experiment is used to close an
infinite quantifier. Literature status is assessed as of 27 November 2026.
\normalsize
\clearpage\tableofcontents\clearpage

\section{The fixed targets and the role of auxiliary questions}
\label{sec:targets}

The requested research pattern is meaningful only if its final step
returns to the original problem:
\[
\text{auxiliary question}\ \longrightarrow\ \text{proved mechanism}
\ \longrightarrow\ \text{original named assertion}.
\]
The first two arrows can be productive without establishing the last
assertion. We keep this distinction explicit throughout.

Our primary target is Legendre's conjecture:
\begin{equation}\label{eq:legendre}
 \forall n\in\Z_{\ge1}\quad\exists p\text{ prime}\quad
 n^2<p<(n+1)^2.
\end{equation}
The 2026 primary discussion of Chamberland and Straub
\cite{CS} retains this assertion as a conjecture. Its conditional
larger-power results do not replace the exponent $2$ in
\eqref{eq:legendre}.

A second target is the Erd\H{o}s--Graham reciprocal-sum rigidity problem,
recorded as Erd\H{o}s Problem~\#243 and restated by Koizumi
\cite[Question 1]{Koizumi}:
\begin{equation}\label{eq:erdos243}
 \begin{gathered}
 1\le a_1<a_2<\cdots,\qquad
 \frac{a_{n+1}}{a_n^2}\longrightarrow1,
 \qquad \sum_{n\ge1}\frac1{a_n}\in\Q
 \\
 \Longrightarrow\quad
 a_{n+1}=a_n^2-a_n+1\quad\text{for all sufficiently large }n.
 \end{gathered}
\end{equation}
The reciprocal ratio convention $a_n^2/a_{n+1}\to1$ is equivalent,
because the terms are positive. This problem is distinct from
Erd\H{o}s Problem~\#263. We neither identify the two problems nor
claim that a result about one automatically resolves the other.

\subsection{What this manuscript adds to the preceding analysis}

The preceding analysis of the constant $D$ established a density bound
for its nonzero digits and excluded algebraic irrational degrees
$2$ through $7$. The present work goes in different directions:

\begin{enumerate}
\item It solves a sharp auxiliary extremal problem about the constant
in a square-root gap-mass bound, including an explicit extremal model.
\item It classifies the possible analytic type of the actual exception
generating function, and gives explicit finite certificates that would
work if their structural hypotheses were proved.
\item It constructs exact obstructions to two sparsity-based transfer
lemmas, and supplies a complete attributed proof for a reciprocal-tail
termination criterion.
\item It tests other named graph problems and records the prior work
that prevents their rediscovery from being counted as a new solution.
\end{enumerate}

None of these statements is presented as a certification that the new
manuscript has greater research significance than every item in the
corpus. Correctness of a lemma, breadth of its applications, historical
novelty, and resolution of a particular conjecture are separate claims.

\subsection{Corpus snapshot and selective synthesis}

The inspected repository \cite{Corpus} had commit
\begin{center}\small
\texttt{e73bae9fd0c319129b19aacc5c317f7abf2865ef}.
\end{center}
The mathematical inputs were selected by relevance; this is not a
certification of all manuscripts in the repository.

\begin{center}\small
\begin{tabularx}{\linewidth}{@{}p{0.12\linewidth}XX@{}}
\toprule
Manuscript & Relevant direction & Consequence for this investigation\\
\midrule
{[173]} & Least Mills constants and large prime gaps & Provides the route
to the external gap-mass estimate used in Section~\ref{sec:baseline}.\\[4pt]
{[013]} & Recognition of bounded Mahler relations & Motivates checking
whether a uniform finite description of the actual exception sequence
has been established. It has not.\\[4pt]
{[232]} & Prime intervals and zeta-zero correlations & Its RH and phase
condition remain hypotheses; positivity alone is not the required
upper estimate.\\[4pt]
{[254]} & Prime-counting exceptions under RH & Larger exponents and
conditional statements do not establish the original square case.\\[4pt]
{[171], [251]} & Harmonic and Randi\'c graph indices & Motivate the graph
tests and an external priority check. Their claims are not imported
as independent established theorems here.\\[4pt]
{[264]} & Graffiti problems and arithmetic graphs & Motivates distinguishing
a graph inequality from a prime-specific conclusion.\\
\bottomrule
\end{tabularx}
\end{center}

\section{The exact exception constant and the available prime input}
\label{sec:baseline}

Write
\[
 c_n=\#\{p\text{ prime}:n^2<p<(n+1)^2\},\quad
 b_n=\floor{\frac1{1+c_n}}=\mathbf1_{\{c_n=0\}},
\]
and define
\[
 \E=\{n\ge1:b_n=1\},\qquad E(N)=\#(\E\cap[1,N]),
 \qquad D=\sum_{n\ge1}b_n3^{-n},\qquad
 F(z)=\sum_{n\ge1}b_nz^n.
\]
The letter $b_n$ is used for the exception mask so that $a_n$ can retain
its conventional role in reciprocal-sum sequences.

\begin{proposition}\label{prop:encoding}
The series $D$ converges, $0\le D\le1/2$, and
\[
 D=0\quad\Longleftrightarrow\quad\E=\varnothing
 \quad\Longleftrightarrow\quad\text{Legendre's conjecture}.
\]
For $D_N=\sum_{n\le N}b_n3^{-n}$,
\[
 0\le D-D_N\le\frac1{2\cdot3^N}.
\]
\end{proposition}
\begin{proof}
Every summand is nonnegative and at most $3^{-n}$. Their geometric
majorant sums to $1/2$, and its tail after $N$ is
$1/(2\cdot3^N)$. A sum of nonnegative terms is zero exactly when all
its terms vanish. The definition of $b_n$ then gives the equivalence.
\end{proof}

Let $p_j$ denote consecutive primes. The analytic input of
J\"arviniemi \cite[Theorem 1.1]{J} is
\begin{equation}\label{eq:j}
 \sum_{\substack{x\le p_j\le2x\\p_{j+1}-p_j\ge\sqrt x}}
 (p_{j+1}-p_j)\ll_\eta x^{57/100+\eta}
 \qquad(\eta>0).
\end{equation}

\begin{proposition}\label{prop:sparse}
For every $\varepsilon>0$,
\begin{equation}\label{eq:sparse}
 E(N)\ll_\varepsilon N^{7/50+\varepsilon}.
\end{equation}
In particular $E(N)=o(N)$.
\end{proposition}
\begin{proof}
Consider empty square intervals indexed by $N\le n<2N$. Each lies
inside a consecutive-prime gap of length at least $2N+1$.
Bertrand's postulate places the preceding prime in
$[N^2/2,4N^2]$ for sufficiently large $N$. Cover this range by the
three dyadic shells with left endpoints $N^2/2,N^2,2N^2$.
The gap qualifies for \eqref{eq:j} in whichever shell contains its
left endpoint. If several empty intervals lie in one prime gap, the
sum of their lengths is at most that gap's length: their interiors
are disjoint. Hence
\[
 (2N+1)\#(\E\cap[N,2N))
 \ll_\eta N^{114/100+2\eta}.
\]
Set $2\eta=\varepsilon$ and sum over dyadic ranges. The finitely many
small indices are absorbed into the implied constant.
\end{proof}

\begin{proposition}\label{prop:rational}
The number $D$ is rational if and only if $\E$ is finite. If
$\E\ne\varnothing$ is finite and $M=\max\E$, its reduced denominator
is exactly $3^M$.
\end{proposition}
\begin{proof}
The ternary expansion of $D$ has digits $b_n\in\{0,1\}$; it is not
an expansion with an eventual tail of $2$'s. A rational number has an
eventually periodic expansion in base $3$. A nonzero eventual period
would give positive asymptotic density to the digit $1$, contradicting
Proposition~\ref{prop:sparse}. Thus rationality forces eventual zero
digits. The converse is immediate. If $M$ is the final nonzero position,
the numerator of $3^MD$ is congruent to $1$ modulo $3$, proving the
denominator assertion.
\end{proof}

For context, the theorem of Adamczewski and Faverjon \cite{AF} gives
at least a constant multiple of $N^{1/d}$ nonzero base-$3$ digits for
an algebraic irrational of degree $d$. Combining it with
\eqref{eq:sparse} excludes $2\le d\le7$, because $7/50<1/7$.
It does not exclude $d=8$, for which $1/8<7/50$. The next sections
test possible ways of extending this information.

\section{A sharp auxiliary gap-mass problem}
\label{sec:threshold}

\begin{definition}
For an infinite set $Q=\{q_1<q_2<\cdots\}$ of positive integers,
put $g_j=q_{j+1}-q_j$ and, for real $x>0$, define
\begin{equation}\label{eq:SQ}
 S_Q(x)=\sum_{\substack{x\le q_j\le2x\\g_j\ge\sqrt x}}g_j.
\end{equation}
The two endpoints in the condition $x\le q_j\le2x$ are included.
\end{definition}

\begin{question}
What is the smallest possible value of
$\limsup_{x\to\infty}S_Q(x)/\sqrt x$ among infinite integer sets
having infinitely many empty intervals $(n^2,(n+1)^2)$?
Does that value increase if one additionally imposes accurate
prime-like counting asymptotics and an arbitrary finite prime prefix?
\end{question}

\begin{theorem}[Sharp auxiliary threshold]\label{thm:sharp}
The answer to the first question is $2\sqrt2$. More precisely:
\begin{enumerate}[label=(\alph*)]
\item Every infinite integer set $Q$ with infinitely many empty
consecutive-square intervals satisfies
\begin{equation}\label{eq:thresholdlower}
 \limsup_{x\to\infty}\frac{S_Q(x)}{\sqrt x}\ge2\sqrt2.
\end{equation}
\item For any prescribed finite prime prefix, equality is attained
by such a set $Q$ satisfying
\begin{align}
 \vartheta_Q(x):=\sum_{\substack{q\in Q\\q\le x}}\log q
 &=x+O(\sqrt x),\label{eq:thetaQ}\\
 Q(x):=\#(Q\cap[1,x])
 &=\int_3^x\frac{dt}{\log t}
   +O\!\left(\frac{\sqrt x}{\log x}\right).
 \label{eq:countQ}
\end{align}
\end{enumerate}
\end{theorem}

\subsection{Proof of the universal lower bound}

Let $(n^2,(n+1)^2)$ be an empty interval sufficiently far out that
$Q$ has an element at most $n^2$. Let $q$ be the last such element
and $q'$ its successor in $Q$. Then
\[
 q\le n^2,\qquad q'\ge(n+1)^2,
 \qquad g=q'-q\ge2n+1+n^2-q>2\sqrt q.
\]
At $x=q/2$, the gap has its left endpoint at the included endpoint
$2x$. It qualifies for \eqref{eq:SQ}, and consequently
\begin{equation}\label{eq:firstgap}
 \frac{S_Q(q/2)}{\sqrt{q/2}}
 \ge\sqrt2\frac g{\sqrt q}>2\sqrt2.
\end{equation}
If there are infinitely many empty square intervals, the corresponding
$q$ are unbounded. Indeed a fixed consecutive gap contains only
finitely many square intervals. Taking a limsup proves part~(a).

\subsection{Construction attaining the constant}

For $x\ge3$, write
\[
 L(x)=\int_3^x\frac{dt}{\log t},\qquad
 B=\{m\in\Z_{\ge4}:\floor{L(m)}>\floor{L(m-1)}\}.
\]
Since $0<L(m)-L(m-1)<1$, every floor increment is $0$ or $1$.
Thus
\begin{equation}\label{eq:baselineB}
 B(x)=\floor{L(\floor{x})}=L(x)+O(1).
\end{equation}
Partial summation yields
\begin{equation}\label{eq:thetaB}
 \vartheta_B(x)=x+O(\log x).
\end{equation}
For completeness, write $B(t)=L(t)+R(t)$ with $R(t)=O(1)$ in
$B(x)\log x-\int_3^x B(t)\,dt/t$. The contribution of $L$ is
$x-3$, and that of $R$ is $O(\log x)$.

The consecutive gaps of $B$ are $O(\log x)$ near $x$.
To see this, choose a fixed sufficiently large $C$.
For all large $x$, the integral of $1/\log t$ over
$(x,x+C\log x]$ exceeds $2$. After allowing for its integer
endpoints, \eqref{eq:baselineB} guarantees an element of $B$ there.

Choose a large integer $n_1$ and define $n_{k+1}=n_k^2$. Delete from
$B$ all its elements in
\[
 I_k=(n_k^2,(n_k+1)^2),\qquad k\ge1,
\]
and call the remaining set $Q$. These intervals are disjoint and
widely separated. By \eqref{eq:thetaB}, the deletion in $I_k$ removes
weighted mass $O(n_k)$. Every removed point has logarithm comparable
to $\log n_k$, so it removes $O(n_k/\log n_k)$ points. Moreover
\[
 \sum_{j\le k}n_j=O(n_k),\qquad
 \sum_{j\le k}\frac{n_j}{\log n_j}
 =O\!\left(\frac{n_k}{\log n_k}\right).
\]
If $I_k$ is the last deletion interval beginning before $x$, then
$n_k\le\sqrt x$. This proves \eqref{eq:thetaQ} and
\eqref{eq:countQ}, including when $x$ lies inside $I_k$; bounding
the partial deletion by the full deletion suffices.

Let $u_k$ be the last point of $B$ at most $n_k^2$, and $v_k$ the first
point at least $(n_k+1)^2$. These are consecutive in $Q$. The gap
estimate for $B$ gives
\begin{equation}\label{eq:mergedgaps}
 u_k=n_k^2+O(\log n_k),\quad u_k\le n_k^2,\qquad
 v_k-u_k=2n_k+1+O(\log n_k).
\end{equation}
All other consecutive $Q$-gaps remain ordinary $B$-gaps. In a
sufficiently large shell $[x,2x]$, at most one $u_k$ occurs, since
$u_{k+1}/u_k\to\infty$. Ordinary gaps have length $O(\log x)$
and do not enter $S_Q(x)$ for large $x$. If the one exceptional gap
does enter, then $x\ge u_k/2$ and
\[
 \frac{S_Q(x)}{\sqrt x}
 \le\sqrt2\frac{v_k-u_k}{\sqrt{u_k}}=2\sqrt2+o(1).
\]
Conversely $x=u_k/2$ makes this gap contribute and the same ratio
tends to $2\sqrt2$. This proves the asserted equality in the limsup.

Finally choose an integer $t\ge2$ with $t^2>M$, and choose $n_1>t$.
Replace $Q\cap[1,M]$ by the set of primes at most $M$, and include
the composite number $t^2$ in $Q$. These changes involve only
finitely many points, so they preserve both asymptotic estimates
and the limiting gap-mass constant. The added square lies before
the first deletion interval, and every $I_k$ remains empty.
This also makes explicit that $Q$ is not the prime set.
Theorem~\ref{thm:sharp} is proved.\hfill$\square$

\subsection{What transfers to the actual primes}

\begin{corollary}\label{cor:strictthreshold}
If the actual prime gaps satisfy
\[
 \limsup_{x\to\infty}\frac{S_{\mathbb P}(x)}{\sqrt x}<2\sqrt2,
\]
then only finitely many consecutive-square intervals are prime-free.
\end{corollary}
\begin{proof}
This is the contrapositive of Theorem~\ref{thm:sharp}(a).
\end{proof}

This corollary is a sufficient condition, not an established estimate
for primes. The available bound \eqref{eq:j} has exponent
$57/100+\eta>1/2$, and supplies no bound of this form.
Even proving the displayed strict limsup would leave a finite
initial range to settle.

There is also an exact pointwise distinction. If
$S_Q(x)\le2\sqrt2\sqrt x$ for every $x\ge X_0$, then
\eqref{eq:firstgap} excludes any empty square interval whose
preceding element is at least $2X_0$. Equality in the \emph{limsup}
does not imply this pointwise inequality. The constructed model has
ratios strictly greater than $2\sqrt2$ at $x=u_k/2$, approaching
$2\sqrt2$ from above.

\statusbox{The extremal set $Q$ is an artificial integer set and contains
composite numbers. Theorem~\ref{thm:sharp} proves the limitation of
an argument based only on these aggregate estimates. It is not a
counterexample to Legendre's conjecture or to the Riemann hypothesis,
and it does not establish the optimal constant for the actual primes.}

\section{Analytic classification of the actual exception series}
\label{sec:function}

\begin{theorem}\label{thm:function}
For the actual function $F(z)=\sum_{n\ge1}b_nz^n$ from
Section~\ref{sec:baseline}, the following are equivalent:
\begin{enumerate}[label=(\roman*)]
\item $\E$ is finite;
\item $F$ is a polynomial;
\item $F$ is rational;
\item $F$ is algebraic over $\C(z)$;
\item $F$ is D-finite, meaning that it satisfies a nonzero homogeneous
linear differential equation with polynomial coefficients.
\end{enumerate}
If $\E$ is infinite, the unit circle is a natural boundary for $F$.
\end{theorem}
\begin{proof}
Finite support and polynomiality are equivalent, and a polynomial
has all the other displayed properties.

Suppose first that $F$ is rational. Beyond a fixed index, its
coefficients satisfy a constant-coefficient linear recurrence. A
state consisting of a fixed number of consecutive coefficients
belongs to a finite set of binary vectors, and the recurrence
determines the next state. Thus the coefficient sequence is
eventually periodic. Proposition~\ref{prop:sparse} forces that
period to consist entirely of zeros. Hence the support is finite.

Now suppose instead that $\E$ is infinite. Because the coefficients
are bounded integers and infinitely many equal $1$, the radius of
convergence is exactly $1$. The P\'olya--Carlson theorem, stated in
\cite{BMW}, says that an integer-coefficient series of radius $1$
is rational or has the unit circle as a natural boundary. The
rational case was just excluded.

An algebraic function has only finitely many possible finite branch
points and poles, as follows from its defining polynomial and
discriminant after removing repeated factors. A D-finite function
can be continued along paths avoiding the finitely many zeros of
the highest nonzero differential coefficient, by the local
existence theorem for linear differential equations. Neither can
have the entire unit circle as a natural boundary. This excludes
both (iv) and (v) when the support is infinite.
\end{proof}

Together with Proposition~\ref{prop:rational}, this also gives
\[
 D\in\Q\quad\Longleftrightarrow\quad F\text{ is D-finite}.
\]
Neither side has been established for the actual $D$. In particular,
proving that $F$ is not D-finite would show that there are infinitely
many exceptions and that $D$ is irrational; no such non-D-finiteness
proof for this particular $F$ is supplied.

The natural-boundary conclusion is a statement about a function.
It is not a proof that the value $F(1/3)$ is transcendental. The
classification uses a classical theorem and the already proved
sparsity estimate; it is not claimed as a historically new method.

\section{Two exact counterexamples to sparsity-based bridges}
\label{sec:sparsebridges}

\subsection{Thin support can have a full eightfold sumset}

One possible continuation of the digit method would try to show that
the support of every polynomial power $F(z)^h$ remains sparse.
The following construction rules out deriving that assertion from
\eqref{eq:sparse} alone.

\begin{proposition}[Classical thin-basis construction]\label{prop:thin}
For every integer $h\ge2$ and cutoff $M\ge0$, there is a set
$A\subset\{M+1,M+2,\ldots\}$ such that
\[
 A(X)=O_h(X^{1/h}),\qquad
 hA=\{h(M+1),h(M+1)+1,\ldots\},
\]
where repetitions are permitted in the $h$-fold sumset.
\end{proposition}
\begin{proof}
For $i=0,\ldots,h-1$, let $B_i$ consist of the nonnegative integers
whose binary $1$-digits occur only in positions congruent to $i$
modulo $h$, and put $B=\bigcup_i B_i$. Splitting the binary digits
of any nonnegative integer according to their positions represents
it as a sum of one member of each $B_i$. Thus $hB=\Z_{\ge0}$.
Among the first $L=\floor{\log_2X}+1$ digit positions, each class
contains at most $\lceil L/h\rceil$ positions, so
\[
 \#(B\cap[0,X])\le h\,2^{\lceil L/h\rceil}=O_h(X^{1/h}).
\]
Take $A=M+1+B$ and translate the sumset identity.
\end{proof}

This is the Raikov--St\"ohr construction, with antecedents recorded
in Nathanson \cite{Nathanson}; it is not a new additive-basis theorem.
At $h=8$, its exponent $1/8$ is smaller than $7/50$. Hence a set
satisfying an even stronger cardinality bound than
\eqref{eq:sparse} can have its eighth sumset contain every
sufficiently large integer. The claim concerns abstract support
data, not the actual prime exceptions, and asserts nothing about
algebraicity of $\sum_{a\in A}3^{-a}$.

\subsection{Extreme sparsity can coexist with maximal block complexity}

\begin{proposition}\label{prop:complexity}
For every cutoff $M$ there is a binary sequence $(u_n)_{n\ge1}$,
zero through $M$, with
\[
 \sum_{n\le X}u_n=O\bigl((\log\log(X+3))^2\bigr),
\]
but with exactly $2^k$ distinct consecutive blocks of length $k$
for every $k\ge1$.
\end{proposition}
\begin{proof}
Enumerate all nonempty finite binary words $w_1,w_2,\ldots$ in
length-lexicographic order. Then $|w_j|\le j$. Insert $w_j$
beginning at position $M+2^{2^j}$, and put zeros at all other
positions. The insertion intervals are disjoint. If the $j$th
insertion has begun by $X$ and the $(j+1)$st has not, at most
$\sum_{i\le j}|w_i|\le j(j+1)/2$ digits equal $1$.
Here $j=O(\log\log(X+3))$. Every binary word occurs as an inserted
block, so all $2^k$ possibilities occur for each length $k$.
\end{proof}

Thus sparsity alone does not supply a low-complexity or finite-state
description of the actual mask. More information about the prime
exceptions would be needed for that route.

\section{Finite certificates with explicit missing hypotheses}
\label{sec:certificates}

A long verified zero prefix becomes a global proof only when a
proved structural rule bounds the position of the first nonzero
coefficient. We give two precise versions.

\begin{theorem}[Differential certificate]\label{thm:ODE}
Let $G\in\C[[z]]$ be nonzero and satisfy
\[
 \sum_{j=0}^r P_j(z)G^{(j)}(z)=0,
\]
where $r\ge1$, the $P_j\in\Z[z]$ are not all zero, and every
coefficient of every $P_j$ has absolute value at most $H\ge1$.
Then
\begin{equation}\label{eq:valuation}
 \ord_0G\le1+H(r+1)r!.
\end{equation}
\end{theorem}
\begin{proof}
Put $m=\ord_0G$. If $m<r$, the assertion is immediate. Otherwise
write $v_j=\ord_0P_j$ and let $\beta_j$ be its lowest nonzero
coefficient, omitting zero polynomials. Set
\[
 \mu=\min_j(v_j-j),\qquad
 I(t)=\sum_{v_j-j=\mu}\beta_j(t)_j,
\]
where $(t)_j=t(t-1)\cdots(t-j+1)$ and $(t)_0=1$.
The coefficient of $z^{m+\mu}$ in the differential identity is
the leading coefficient of $G$ multiplied by $I(m)$; hence $I(m)=0$.
The polynomial $I$ is nonzero, since its summands have distinct
degrees and the leading coefficient of the highest-degree summand
cannot cancel. Its coefficients are integers.

The sum of the absolute values of the coefficients of $(t)_j$
is $j!$. Consequently every coefficient of $I$ has absolute value
at most $H\sum_{j=0}^rj!\le H(r+1)r!$. Its leading coefficient
has absolute value at least $1$. Cauchy's root bound now gives
\eqref{eq:valuation}.
\end{proof}

If an identity of this type were proved for the actual $F$, checking
zero coefficients through the displayed bound would prove $F=0$.
No such identity with controlled $r,H$ has been obtained. An identity
fitted to a finite zero prefix does not prove that it holds for the
infinite series.

\begin{proposition}[Automaton certificate]\label{prop:automaton}
Suppose a binary sequence is produced from the canonical base-$b$
representations of positive integers by a deterministic finite
automaton with output and at most $s$ states. If some term is $1$,
one such term occurs at an index less than $b^s$.
\end{proposition}
\begin{proof}
Choose a positive input word ending in a state with output $1$.
Keep its first, nonzero digit fixed. From the resulting state there
is a shortest path to a suitable output state, and a shortest path
has no repeated states. It has at most $s-1$ edges. Thus there is
an accepted positive representation with at most $s$ digits,
representing an integer less than $b^s$.
\end{proof}

Sorenson and Webster's verification \cite{SW} covers
$n\le7.05\cdot10^{13}$, which exceeds
$2^{46}=70\,368\,744\,177\,664$. A proved binary presentation of
the actual exception mask with at most $46$ states would therefore
be sufficient. Its existence is not proved by the computation or
by Proposition~\ref{prop:automaton}. These are conditional uses
of completely proved certificate lemmas, not completed certificates
for Legendre's conjecture.

\section{Rational reciprocal tails and an arithmetic barrier}
\label{sec:reciprocal}

This section returns to the original assertion \eqref{eq:erdos243}.
The reduction is the pseudo-greedy framework of Koizumi \cite{Koizumi}.
The one-sided bounded-defect theorem below has an earlier counterpart
in Cook \cite{Cook}; Bado \cite{Bado} gives the two-sided bounded-gap
result. We give an independently checked proof to identify precisely
what it establishes and what it leaves open.

\subsection{Exact variables, with no change in the original assertion}

Assume the hypotheses on the first two lines of \eqref{eq:erdos243},
and write $r=\sum_n1/a_n=p/q\in\Q$ with $p,q>0$. Define
\begin{align*}
 x_n&=r-\sum_{j<n}\frac1{a_j}=\sum_{j\ge n}\frac1{a_j},\\
 d_n&=q\prod_{j<n}a_j,\qquad c_n=d_nx_n\in\Z_{>0},\\
 e_n&=d_n-(a_n-1)c_n,\qquad \varepsilon_n=e_n/c_n.
\end{align*}
The numerators $c_n$ are not reduced; this convention is essential.
Multiplication of the initial fraction by a fixed positive integer
only scales the integer variables and does not change the existence
of a fixed eventual bound on the negative defects.

\begin{lemma}\label{lem:orbit}
These variables satisfy
\begin{equation}\label{eq:orbit}
 c_{n+1}=c_n-e_n,\qquad d_{n+1}=a_nd_n,
 \qquad c_na_n=d_n+c_{n+1},
\end{equation}
and $\varepsilon_n\to0$. Eventually $a_n$ is the unique integer
nearest $x_n^{-1}+1$. Once $e_n=0$ at a sufficiently late index,
all later defects are zero and the Sylvester recurrence follows.
\end{lemma}
\begin{proof}
The tail identity $x_{n+1}=x_n-1/a_n$ gives all three exact
equalities in \eqref{eq:orbit}. Approximate squaring gives
\[
 x_n=\frac1{a_n}+\frac{1+o(1)}{a_n^2}+O(a_n^{-4}),
 \qquad x_n^{-1}=a_n-1+o(1).
\]
For the remainder estimate, use $a_{n+1}\ge a_n^2/2$ eventually
and bound the tail from $a_{n+2}$ onward by a constant multiple
of $1/a_{n+2}$. Thus
$\varepsilon_n=x_n^{-1}+1-a_n\to0$, and eventual nearest-integer
rounding follows.

If $e_n=0$, then $x_n=1/(a_n-1)$, so
$x_{n+1}=1/[a_n(a_n-1)]$. At sufficiently late indices the
nearest-integer rule forces
$a_{n+1}=a_n(a_n-1)+1$ and $e_{n+1}=0$. Iterate.
\end{proof}

\subsection{The bounded-negative-defect theorem, with attribution}

\begin{theorem}[Bounded negative defects; cf.\ Cook]\label{thm:CRT}
Under the hypotheses of \eqref{eq:erdos243}, suppose additionally
that for some fixed integer $K\ge0$,
\begin{equation}\label{eq:boundedneg}
 e_n\ge-K\qquad\text{for every sufficiently large }n.
\end{equation}
Then $e_n=0$ eventually, and
$a_{n+1}=a_n^2-a_n+1$ eventually.
\end{theorem}
\begin{proof}
Discard a finite prefix so that \eqref{eq:boundedneg} and the
rounding conclusion of Lemma~\ref{lem:orbit} hold thereafter.
If $K=0$, then $c_{n+1}\le c_n$. A nonincreasing sequence of positive
integers eventually stabilizes, forcing $e_n=0$.

Let $K\ge1$, and suppose for contradiction that there is no zero
defect in this tail. As $e_n$ is a nonzero integer,
\begin{equation}\label{eq:cdiverges}
 \frac1{c_n}\le\frac{|e_n|}{c_n}\longrightarrow0,
 \qquad\text{so }c_n\longrightarrow\infty.
\end{equation}
On the other hand $c_{n+1}\le c_n+K$, and hence $c_n\le C+Kn$
for a constant $C$.

\emph{Infinitely many denominator primes.}
For $i<j$, $a_i\mid d_j$, and \eqref{eq:orbit} implies
\begin{equation}\label{eq:smallgcd}
 \gcd(a_i,a_j)\mid c_{j+1}.
\end{equation}
Suppose the prime divisors of all sufficiently late $a_n$ belonged
to a finite set $S$ of $s$ primes. In a window of $s+1$ consecutive
late indices $n,\ldots,n+s$, let $T=C+K(n+s+1)$. Approximate
squaring implies growth faster than any polynomial, so every
$a_i$ in this window exceeds $T^s$. Each therefore contains a
prime-power factor $p^{v_p(a_i)}>T$ with $p\in S$. The same prime
cannot serve two indices: their gcd would exceed $T$, contrary
to \eqref{eq:smallgcd}. This would require $s+1$ distinct primes
from a set of $s$, a contradiction.

Choose distinct primes $p_1,\ldots,p_K>K$ which all divide $d_N$
at some sufficiently late common time $N$. They divide every
subsequent $d_n$, by the multiplicative denominator update.

\emph{Persistent divisibility forbids numerator values.}
If $p_i\mid c_n$ at any time $n\ge N$, then also $p_i\mid e_n$.
The recurrences preserve $p_i\mid c_m,d_m$ for every later $m$.
All such $e_m$ are multiples of $p_i$, and $e_m\ge-K>-p_i$,
so they are nonnegative. Under the nonzero-defect assumption
they are at least $p_i$, forcing the positive integers $c_m$
to decrease by at least $p_i$ at every step. This is impossible.
Thus the nonterminating tail can visit no multiple of any $p_i$.

\emph{A finite forbidden block.}
By the Chinese remainder theorem, choose an integer $X>c_N$ with
\[
 X+j\equiv0\pmod{p_{j+1}}\qquad(0\le j<K).
\]
All $K$ consecutive integers $X,\ldots,X+K-1$ are forbidden to
the later numerator sequence. Yet \eqref{eq:cdiverges} says it
eventually crosses $X$. At its first crossing it starts at a value
at most $X-1$ and moves upward by at most $K$, so it lands in that
forbidden block. This contradiction proves eventual zero defects.
Lemma~\ref{lem:orbit} supplies the original recurrence.
\end{proof}

The argument uses an actual arithmetic restriction: prime divisors
persist in $d_n$, while sufficiently large shared divisors force
descent of $c_n$. It is more than a change of coordinates. However,
the same theorem and the same essential CRT-barrier mechanism
already appear in the prior source \cite{Cook}. It therefore does
not meet the requested novelty condition as a new discovery here.

\subsection{A rate condition and the remaining unbounded case}

Put
\[
 \gamma_n=\frac{a_n^2}{a_{n+1}}-1.
\]
The following useful specialization is likewise covered by the
earlier reciprocal-tail work \cite{Cook}.

\begin{corollary}\label{cor:rate}
Under the original rationality and growth assumptions, suppose
that for constants $A\ge0$ and $\delta>0$,
\begin{equation}\label{eq:rate}
 \gamma_n\le\frac1n+\frac{A}{n^{1+\delta}}
 \qquad(n\text{ sufficiently large}).
\end{equation}
Then the Sylvester recurrence holds eventually.
\end{corollary}
\begin{proof}
The exact tail relations give
\[
 a_{n+1}=\frac{a_n^2}{1-\varepsilon_n}-a_n+1-\varepsilon_{n+1}.
\]
As $\varepsilon_n\to0$, division and inversion show
\[
 \gamma_n=-\varepsilon_n+O(1/a_n).
\]
Writing $\varepsilon_n^- =\max(-\varepsilon_n,0)$, we obtain
\[
 \varepsilon_n^-\le\frac1n+\frac{A}{n^{1+\delta}}+O(1/a_n).
\]
Approximate squaring makes $\sum 1/a_n$ convergent rapidly. Since
$c_{n+1}=c_n(1-\varepsilon_n)\le c_n(1+\varepsilon_n^-)$,
taking logarithms and summing gives $c_n=O(n)$. Therefore
$(-e_n)_+=c_n\varepsilon_n^-=O(1)$, and
Theorem~\ref{thm:CRT} applies.
\end{proof}

The original problem assumes only $\gamma_n\to0$, not
\eqref{eq:rate}. It supplies only $e_n=o(c_n)$, not the fixed
lower bound \eqref{eq:boundedneg}. In any possible nonterminating
counterexample, the negative defects must consequently have
unbounded magnitude, even though their ratio to $c_n$ tends to
zero. The first-crossing proof does not prevent jumps over every
fixed forbidden block when the allowed upward jump sizes are
unbounded. This is the remaining step, and it is not proved here.

\section{Exact graph tests and the prior-resolution filter}
\label{sec:graphs}

The corpus suggested looking for a named graph conjecture whose
structure might permit a complete argument. Several concrete
calculations succeeded, but external sources prevented them from
being counted as a previously unavailable resolution.

\subsection{A harmonic-index family that is not a new disproof}

For a simple graph, define its harmonic index by
\[
 H(G)=\sum_{uv\in E(G)}\frac2{d_u+d_v}.
\]
Let $S_{n,k}=K_k\vee\overline K_{n-k}$ be the complete split
graph. Cheng and Wang \cite{CW} conjectured
$H(G)\ge H(S_{n,k})$ for the stated range
$\delta(G)\ge k$, $1\le k\le\lfloor n/2\rfloor+1$.

\begin{proposition}\label{prop:harmonic}
For each integer $t\ge2$, let $G_t=K_{2t}\vee(tK_2)$. It has
$n=4t$, $\delta=2t+1$, and
\[
 H(S_{4t,2t+1})-H(G_t)
 =\frac{(t-1)^2}{3t(4t-1)(2t+1)}>0.
\]
\end{proposition}
\begin{proof}
The $2t$ clique vertices have degree $4t-1$, and the $2t$ matching
vertices have degree $2t+1$. Separating clique, matching, and
cross edges gives
\[
 H(G_t)=\frac{t(2t-1)}{4t-1}+\frac{t}{2t+1}+\frac{4t}{3}.
\]
For the split graph,
\[
 H(S_{4t,2t+1})=\frac{t(2t+1)}{4t-1}
                  +\frac{4t^2-1}{3t}.
\]
Subtracting and putting the terms over a common denominator yields
the claimed expression.
\end{proof}

At $t=2$ the difference is $1/210$. The family lies within the
original boundary range; it is not a change of problem. But Ali
\cite{Ali} had already disproved the conjecture in 2018, and
Sun, Gao, Du, and Xu \cite{Sun} both establish the range $k\le n/2$
and give boundary counterexamples. This calculation is therefore
an exact reconstruction, not a new resolution of an open problem.

\subsection{A proposed stronger distance lemma is false}

Let $\mu(G)$ be the average distance over unordered pairs,
$\alpha(G)$ the independence number, and $\operatorname{diam}(G)$
the diameter. One proposed auxiliary bound was
\[
 2\mu(G)\le\max\{\alpha(G),\operatorname{diam}(G)+1\}.
\]
Take a graph on vertices $0,1,2,3,4,5$, with $0$ adjacent to every
other vertex and additional edges $23,45$. It has seven edges,
diameter $2$, and independence number $3$. Of its $15$ unordered
pairs, seven have distance $1$ and eight have distance $2$.
Consequently
\[
 2\mu(G)=\frac{46}{15}>3
 =\max\{\alpha(G),\operatorname{diam}(G)+1\}.
\]
This is a counterexample to the proposed auxiliary lemma, not to
the induced-linear-forest conjecture. That named conjecture was
posed by Hansen et al.\ \cite{HHKMS}; Bouchard and Hertz
\cite{BH} posted a full-proof claim on 9 November 2026. A focused
reading of a preprint does not certify its entire proof, but the
public prior claim must be considered before asserting originality.

\subsection{A broad graph-energy family supports the conjectured bound}

Write $\mathcal{E}_{\rm gr}(G)$ for adjacency energy, the sum of the
absolute values of the adjacency eigenvalues. A tested candidate
for violating $\mathcal{E}_{\rm gr}(G)\ge\Delta(G)+\delta(G)$ was
a clique joined to disjoint small cliques.

\begin{proposition}\label{prop:energy}
If
\[
 G=K_s\vee\left(\bigsqcup_{i=1}^tK_{r_i}\right),\qquad
 s\ge1,\quad t\ge2,\quad r_i\ge2,
\]
then $G$ is nonsingular and
$\mathcal{E}_{\rm gr}(G)>\Delta(G)+\delta(G)$.
\end{proposition}
\begin{proof}
Put $n=s+\sum_i r_i$ and $r_* =\min_i r_i$.
There are $n-t-1$ independent eigenvectors of eigenvalue $-1$,
obtained by taking vectors of sum zero within each clique cell.
The remaining eigenvalues are those of the symmetric quotient
\[
 Q=\begin{pmatrix}
 s-1&\sqrt{sr_1}&\cdots&\sqrt{sr_t}\\
 \sqrt{sr_1}&r_1-1&&0\\
 \vdots&&\ddots&\\
 \sqrt{sr_t}&0&&r_t-1
 \end{pmatrix}.
\]
Its lower diagonal block is positive definite, while its Schur
complement is
$s-1-s\sum_i r_i/(r_i-1)<0$. Thus $Q$ has exactly one negative
eigenvalue $\lambda$ and $t$ positive eigenvalues, with no zero.
The graph contains an induced three-vertex path; interlacing
gives its least eigenvalue at most $-\sqrt2$. Therefore
$\lambda\le-\sqrt2<-1$.

The adjacency trace is zero, so
$\mathcal{E}_{\rm gr}(G)=2(n-t-1-\lambda)$.
Also $\Delta=n-1$ and $\delta=s+r_*-1$. Choosing an index $i_*$
with $r_{i_*}=r_*$, we obtain
\[
 \mathcal{E}_{\rm gr}(G)-\Delta-\delta
 =\sum_i(r_i-2)-r_*-2\lambda
 >\sum_{i\ne i_*}(r_i-2)\ge0.
\]
\end{proof}

This eliminates a broad proposed counterexample family. It is not
a new proof of the general named conjecture. Mojallal's August
2026 preprint \cite{Mojallal} already claims a more general energy
bound and explicitly lists the $\Delta+\delta$ inequality among
its consequences. We do not promote that public claim to a
fully independently verified theorem, and we do not claim priority
over it.

\section{Assessment against the six required conditions}
\label{sec:gates}

The gates apply to a resolution of an original named problem.
They cannot be passed by relabelling an auxiliary theorem as that
resolution.

\begin{longtable}{@{}p{0.08\linewidth}>{\raggedright\arraybackslash}p{0.26\linewidth}>{\raggedright\arraybackslash}p{0.60\linewidth}@{}}
\toprule
Gate & Requirement & Outcome\\
\midrule\endhead
G1 & Pre-existing named problem & Legendre and the Erd\H{o}s--Graham
reciprocal-sum problem are genuine fixed targets. The sharp
$2\sqrt2$ auxiliary problem is explicitly not substituted for them.\\[7pt]
G2 & Original assertion retained & The original square intervals
and reciprocal recurrence are kept. All transformations used in
the arithmetic arguments are derived by exact equalities. Added
hypotheses are displayed as added hypotheses.\\[7pt]
G3 & Genuinely unresolved at the time & Checked primary sources
retain the two arithmetic targets as unresolved. The harmonic
counterexample is already known; other graph candidates have
public prior complete-proof claims. No exhaustive status or
priority certification is asserted.\\[7pt]
G4 & Unconditional complete resolution & \textbf{Not achieved for
the original targets.} The auxiliary theorems are fully proved,
but no all-index prime-existence estimate and no unrestricted
termination theorem is obtained.\\[7pt]
G5 & Substantive new method resolving the target & \textbf{Not
established.} The sharp model identifies a precise limitation.
The CRT theorem has real arithmetic content, but its needed
bounded-jump hypothesis has not been derived from the original
assumptions, and its mechanism has prior antecedents.\\[7pt]
G6 & No prior identical solution or essential proof & \textbf{Not
certified.} Exact prior counterparts were found for the harmonic
and bounded-defect results. No firstness claim is made for the
remaining auxiliary constructions or elementary consequences.\\
\bottomrule
\end{longtable}

\subsection{The exact remaining transfers}

For Legendre, a conclusion about every original interval remains
necessary. One precise formulation is
\[
 \vartheta((n+1)^2)-\vartheta(n^2)>0\quad(n\ge1),
\]
where $\vartheta$ sums over primes, not prime powers. The square
endpoints are not primes, so this positivity has exactly the
required meaning. Theorem~\ref{thm:sharp} identifies a sharp
limitation of certain aggregate estimates; it does not prove
this positivity. A finite exception set is still different from
an empty exception set.

For the reciprocal-sum target, the exact unclosed implication is
\[
 \text{the integer tail dynamics and }e_n/c_n\to0
 \quad\Longrightarrow\quad e_n=0\text{ eventually}.
\]
The added bound $e_n\ge-K$ closes a subclass. Neither its removal
nor a replacement controlling unbounded upward jumps is supplied.

New auxiliary questions are useful when their solutions supply
information that the original object is proved to possess.
Here several auxiliary questions have complete answers, including
a sharp constant and explicit countermodels. The missing transfer
statements remain explicit rather than being replaced by numerical
evidence, equivalent encodings, or a priority claim based on a
failed search.

\appendix
\section{Verification and provenance notes}

The displayed harmonic difference was checked by exact rational
algebra, including $t=2$, where
$H(G_2)=412/105$ and $H(S_{8,5})=55/14$.
The six-vertex distance calculation uses all $\binom62=15$ pairs.
The sharp threshold argument and the reciprocal-tail CRT proof
received separate independent checks during preparation.
These checks are mathematical review aids, not claims of formal
proof-assistant certification.

Cook's manuscript source was located in the repository
\texttt{wcook04/plectis-erdos}. Its observed remote HEAD was
\begin{center}\small
\texttt{a501de956765730ba4e6c9ad284e6d65855d533a}.
\end{center}
The source inspected contains the theorem named ``bounded negative
part'' and is dated September 2026. Bado's separate preprint is
credited to Idriss Olivier Bado, not to Cook. The public dates and
the mathematical overlap are enough to exclude a first-discovery
claim for the reconstructed result; no blanket certification of
the authors' other claims is implied.

The corpus remains an input of research drafts. The proof of any
result used as an established external input is cited separately.
The named-problem assessment is dated and selective, and the
failure of G4--G6 is stated in the abstract as well as in
Section~\ref{sec:gates}.

To reproduce the PDF in a standard TeX installation, run
\texttt{pdflatex} twice on the supplied \texttt{.tex} file so that
the contents, citations, and cross-references resolve.

\small
\begin{thebibliography}{99}

\bibitem{Corpus}
QEEDDD, \emph{Verifiable Mathematical Reasoning Corpus}, research
repository, snapshot \texttt{e73bae9fd0c319129b19aacc5c317f7abf2865ef},
inspected 27 November 2026.
\url{https://github.com/QEEDDD/Verifiable-Mathematical-Reasoning-Corpus}.
Manuscript numbers in Section~\ref{sec:targets} refer to this snapshot.

\bibitem{J}
O. J\"arviniemi, \emph{On large differences between consecutive primes},
arXiv:2212.10965v2 (2022), Theorem 1.1.
\url{https://arxiv.org/abs/2212.10965v2}.

\bibitem{AF}
B. Adamczewski and C. Faverjon,
\emph{Chiffres non nuls dans le d\'eveloppement en base enti\`ere des
nombres alg\'ebriques irrationnels},
C. R. Acad. Sci. Paris, S\'er. I \textbf{350} (2012), 1--4.
\href{https://doi.org/10.1016/j.crma.2011.12.002}{doi:10.1016/j.crma.2011.12.002}.

\bibitem{SW}
J. Sorenson and J. Webster, \emph{An algorithm and computation to verify
Legendre's conjecture up to $7\cdot10^{13}$}, Research in Number
Theory \textbf{11} (2025), article 4; online 6 December 2024.
\href{https://doi.org/10.1007/s40993-024-00589-4}{doi:10.1007/s40993-024-00589-4}.

\bibitem{CS}
M. Chamberland and A. Straub, \emph{Weakening the Legendre Conjecture},
American Mathematical Monthly \textbf{133} (2026), 855--865.
\href{https://doi.org/10.1080/00029890.2026.2694309}{doi:10.1080/00029890.2026.2694309}.
\href{https://arminstraub.com/downloads/pub/legendre-riemann.pdf}{Author preprint}.

\bibitem{BMW}
J. Bell, R. Miles, and T. Ward, \emph{Towards a P\'olya--Carlson
dichotomy for algebraic dynamics}, Indagationes Mathematicae
\textbf{25} (2014), 652--668.
\url{https://arxiv.org/abs/1307.2369}.

\bibitem{Nathanson}
M. B. Nathanson, \emph{Waring's problem for finite intervals},
Proceedings of the American Mathematical Society \textbf{96}
(1986), 15--17. The introduction records the thin-basis
constructions of Raikov and St\"ohr (1937).
\href{https://www.theoryofnumbers.com/melnathanson/pdfs/nath1986-53.pdf}{Author-hosted article}.

\bibitem{Koizumi}
J. Koizumi, \emph{Irrationality of the reciprocal sum of doubly
exponential sequences}, Integers \textbf{26} (2026), A28,
published 20 February 2026.
\href{https://doi.org/10.5281/zenodo.18714404}{doi:10.5281/zenodo.18714404}.
\url{https://math.colgate.edu/~integers/aa28/aa28.pdf}.
Question 1, Corollary 3, and Lemmas 2--4.

\bibitem{Bado}
I. O. Bado, \emph{Prime-Support Rigidity and Primitive Pseudo-Greedy
Dynamics: Partial Progress on Erd\H{o}s Problem \#243}, public
preprint, September 2026, Theorem 5.1.
\href{https://doi.org/10.13140/RG.2.2.36612.08325}{doi:10.13140/RG.2.2.36612.08325}.

\bibitem{Cook}
W. Cook, \emph{Reciprocal Sums and the Sylvester Recurrence},
public working paper, September 2026; theorem named
``bounded negative part'' and associated rate corollaries.
\href{https://github.com/wcook04/plectis-erdos/blob/a501de956765730ba4e6c9ad284e6d65855d533a/paper/243/erdos-243-reciprocal-tail-rigidity.tex}{Pinned manuscript source}.
Accessed 27 November 2026. This citation does not certify all
claims in the associated repository.

\bibitem{CW}
M. Cheng and L. Wang, \emph{A lower bound for the harmonic index
of a graph with minimum degree at least three}, Filomat
\textbf{30}(8) (2016), 2249--2260.
\href{https://doi.org/10.2298/FIL1608249C}{doi:10.2298/FIL1608249C}.

\bibitem{Ali}
A. Ali, \emph{Counter examples to a conjecture concerning harmonic
index}, Asian-European Journal of Mathematics \textbf{11}(3)
(2018), 1850035.
\href{https://doi.org/10.1142/S1793557118500353}{doi:10.1142/S1793557118500353}.

\bibitem{Sun}
X. Sun, Y. Gao, J. Du, and L. Xu, \emph{On a conjecture of the
harmonic index and the minimum degree of graphs}, Filomat
\textbf{32}(10) (2018), 3435--3441.
\href{https://doi.org/10.2298/FIL1810435S}{doi:10.2298/FIL1810435S}.

\bibitem{HHKMS}
P. Hansen, A. Hertz, R. Kilani, O. Marcotte, and D. Schindl,
\emph{Average distance and maximum induced forest}, Journal
of Graph Theory \textbf{60}(1) (2009), 31--54; first published
online 30 July 2008.
\href{https://doi.org/10.1002/jgt.20344}{doi:10.1002/jgt.20344}.

\bibitem{BH}
M. Bouchard and A. Hertz, \emph{Average distance and maximum
induced linear forest}, arXiv:2611.12781v1, 9 November 2026.
\url{https://arxiv.org/abs/2611.12781}.
Public full-proof claim; not independently certified in full here.

\bibitem{Mojallal}
S. A. Mojallal, \emph{Rank--Average-Degree Bound for Graph Energy},
arXiv:2608.22139v2, August 2026.
\url{https://arxiv.org/abs/2608.22139}.
Public full-proof claim; not independently certified in full here.

\end{thebibliography}
\end{document}
